\label{sec:methodology}

\subsection{Section~2 as the Primary Post-\textit{Shelby} Tool}

Section~2 of the Voting Rights Act, 52~U.S.C.~\S~10301, prohibits any
voting qualification, standard, practice, or procedure that ``results in''
the denial or abridgment of the right to vote on account of race or color.
As amended in 1982, Section~2 employs a ``results'' standard: plaintiffs
need not prove discriminatory intent but only that the challenged practice
has a discriminatory effect. In the redistricting context, Section~2 is
primarily enforced through the vote-dilution framework of
\textit{Thornburg v.\ Gingles}.\footnote{\textit{Thornburg v.\ Gingles},
478 U.S.\ 30 (1986).}

\textit{Gingles} established three preconditions for a Section~2 claim
that a redistricting plan unlawfully dilutes minority voting strength
by failing to create a majority-minority district:

\begin{enumerate}
\item \textit{Numerosity and compactness}: The minority group is ``sufficiently
  large and geographically compact to constitute a majority in a
  single-member district.''

\item \textit{Political cohesion}: The minority group is ``politically
  cohesive''---members of the group tend to vote together.

\item \textit{Majority bloc voting}: The white majority votes sufficiently
  as a bloc to defeat the minority group's preferred candidates in the
  absence of a majority-minority district.
\end{enumerate}

If all three preconditions are met, the court examines the ``totality of
circumstances'' to determine whether the challenged plan ``results in''
minority vote dilution. The analysis includes, but is not limited to,
the history of discrimination, the extent of racially polarized voting,
the use of unusually large election districts, and the lack of success
of minority candidates in the jurisdiction.

\subsection{\textit{Abbott v.\ Perez} (2018) and the Intent Bar}

\textit{Abbott v.\ Perez}\footnote{\textit{Abbott v.\ Perez}, 585 U.S.\ 579
(2018).} addressed the Texas redistricting litigation that had been running
since the 2012 cycle. The Supreme Court, in an opinion by Justice Alito,
reversed the district court's findings that the majority of challenged
Texas congressional and state legislative districts violated the VRA or the
Equal Protection Clause. The Court imposed a high burden on plaintiffs
seeking to prove discriminatory intent, requiring evidence sufficient to
overcome the presumption of legislative good faith.

For Section~2 redistricting claims based on discriminatory intent (as
opposed to the results-based dilution analysis of \textit{Gingles}),
\textit{Abbott v.\ Perez} makes it substantially harder to prevail. But
the decision did not alter the \textit{Gingles} results-based framework;
it addressed the intent prong that applies when plaintiffs allege racial
gerrymandering under the Equal Protection Clause rather than vote dilution
under Section~2.

\subsection{\textit{Brnovich v.\ Democratic National Committee} (2021)}

\textit{Brnovich v.\ Democratic National Committee}\footnote{\textit{Brnovich
v.\ Democratic National Committee}, 594 U.S.\ 647 (2021).} is the most
significant post-\textit{Shelby} weakening of Section~2, but its scope
is limited to vote-denial claims. The case involved two Arizona election
rules: a policy of not counting out-of-precinct provisional ballots and
a law prohibiting ballot collection by third parties. The five-Justice
majority, per Justice Alito, held that these rules did not violate
Section~2's results test, and in doing so articulated a multi-factor
framework for evaluating non-redistricting Section~2 claims.

The \textit{Brnovich} framework applies to Section~2 vote-denial claims
(challenges to election rules that burden minority voters' ability to vote)
rather than Section~2 vote-dilution claims (challenges to redistricting
plans that dilute minority voting strength). Justice Kagan's dissent,
joined by Justices Breyer and Sotomayor, argued that the majority had
effectively created a new test for vote-denial claims that was less
demanding of defendants than the statutory text required.

For redistricting practitioners, \textit{Brnovich}'s key significance
is what it \textit{did not} do: it expressly preserved the Section~2
vote-dilution framework for redistricting claims, and it did not alter
the \textit{Gingles} analysis. The Court's opinion noted that the
two categories of Section~2 claim---vote denial and vote dilution---
are analyzed differently, and that the \textit{Brnovich} analysis was
directed at vote-denial claims only.\footnote{\textit{Brnovich},
594 U.S.\ at 680 n.19.} The \textit{Gingles} framework thus survived
\textit{Brnovich} intact for redistricting purposes.

\subsection{Congressional Inaction: The Coverage Gap Persists}

The John Lewis Voting Rights Advancement Act (H.R.~4), named for the late
Georgia congressman and civil rights icon, would establish a new coverage
formula for Section~5 based on a different methodology: jurisdictions
with more than a threshold number of voting rights violations in the
preceding twenty-five years would be covered. The Act passed the House
in December 2019 and again in August 2021 but has not received a Senate
vote as of 2026.

The failure to restore Section~5 preclearance means that the 2020 and
2030 redistricting cycles in former covered states proceed entirely under
Section~2, without the ex ante review that Section~5 provided. Former
covered states retain no federal obligation to submit redistricting plans
for review; minority plaintiffs must file suit after enactment and litigate
for years before courts may order remedial maps.

This asymmetry---ex ante administrative review under Section~5 versus
ex post judicial review under Section~2---is the defining structural
feature of the post-\textit{Shelby} redistricting environment for the
foreseeable future.
